% problem-000349 3 多钒酸盐组成分析
\section*{3. 多钒酸盐组成分析}
实验1：准确称取0.9291 g A于三颈瓶中，加100 ml 蒸馏⽔和150 mL 20% NaOH溶液，加热煮沸，⽣成的氨⽓⽤50.00 mL 0. 1000 mol L−1 HCl 标准溶液吸收。加⼊酸碱指⽰剂，⽤0.1000 mol L−1 $\mathrm { N a O H }$ 标准溶液滴定剩余的HCl标准溶液，终点时消耗19.88 mL NaOH 标准溶液。

实验 2：确称取 0.3097 g A 于锥形瓶中，加⼊40 mL 1.5 mol L−1 H ${ } _ { 2 } S O _ { 4 }$ ，微热使之溶解。加⼊50 mL蒸馏⽔和 1 $\mathrm { g N a H S O _ { 3 } , }$ 搅拌5分钟，使反应完全，五价钒被还原成四价。加热煮沸15分钟，然后⽤0.02005 mol L−1$\mathrm { K M n O } _ { 4 }$ 标准溶液滴定，终点时消耗25.10 mL $\mathrm { K M n O } _ { 4 }$ 标准溶液。

\subsection*{3-1}
图3.1为滴定曲线图，请回答哪⼀个图为实验1的滴定曲线；请根据此滴定曲线选择⼀种最佳酸碱指⽰剂，并简述做出选择的理由。有关指⽰剂的变⾊范围如下：甲基橙(pH 3.1\~4.4)，溴甲酚绿(pH3.8\~5.4)，酚酞(pH 8.0\~10.0)。

（图：）
(a)

（图：）
(b)
图 3.1

（图：）
(c)

\subsection*{3-2}
在实验2中，加 $\mathrm { N a H S O _ { 3 } }$ 还原 $\mathrm { V _ { 1 0 } O _ { 2 8 } } ^ { 6 - }$ ，反应完全后需要加热煮沸15分钟。煮沸的⽬的是什么？写出$\mathrm { K M n O } _ { 4 }$ 滴定反应的离⼦⽅程式。

\subsection*{3-3}
根据实验结果，计算试样 A 中 $\mathrm { N H } _ { 4 }$ +和 $\mathrm { V _ { 1 0 } O _ { 2 8 } } ^ { 6 - }$ −的质量分数，结合元素分析结果确定A的化学式（�和�取整数）。
